% ---------------------------------------------------------------------------
% Introduction
% Drafted from experiments/Overview.md and experiments/Report.md.
% Every number here is reproduced by experiments/verify_all.py (214 checks).
% Citation keys are placeholders -- replace with the keys in the .bib.
% ---------------------------------------------------------------------------

\section{Introduction}

Vision-language-action (VLA) policies are slow. On the hardware we used, a
single observation costs one of our backbones 2.8~seconds. This has motivated
a growing set of \emph{training-free} interventions that leave the checkpoint
untouched and make inference cheaper: skipping decoder layers, caching or
pruning visual tokens, and reusing a predicted action across several control
steps~\cite{efficientvla, vlacache, fastv, molevla, shortgpt}.

Papers proposing such interventions report their result in a recurring form:

\begin{quote}
\emph{Applying method $M$ to backbone $A$ on benchmark $B$ reduced FLOPs by
$X\%$ while success dropped only $Y$ points. Therefore $M$ is efficient.}
\end{quote}

For that final clause to follow, the measured effect has to be a property of
\emph{the method} rather than of the single configuration it was measured in.
Changing the backbone, the benchmark, or an implementation detail that the
paper does not report should not change the direction of the result. That
premise, rather than any individual method, is what this paper tests.

We re-measure three existing interventions under one protocol across three
open VLA backbones (OpenVLA~\cite{openvla}, SpatialVLA~\cite{spatialvla},
UniVLA~\cite{univla}) and two SimplerEnv~\cite{simplerenv} suites
(WidowX-Bridge and Google Robot / Fractal). The three interventions were
chosen to consume different resources: \emph{action repeat} changes how often
the policy is called, \emph{foveation} changes what the policy is shown, and
\emph{depth pruning} changes how much of the network each call uses. The grid
has five filled cells (UniVLA releases a Bridge-only checkpoint) and eight
conditions per cell, for 7{,}198 episodes.

The measurement is paired. We do not subtract two aggregate success rates:
every condition replays the same initial states as its cell's baseline, and we
count only the episodes whose outcome changed, testing them with McNemar's
exact test. The simulator is deterministic --- we verified that re-running a
condition reproduces it episode for episode --- so there is no run-to-run
variance, and the remaining uncertainty is entirely in the $p$-values. Across
the 38 paired tests in the grid, a Bonferroni threshold of
$\alpha \approx 0.0013$ leaves eight cells significant; across the 43
cross-cell comparisons, seven survive $\alpha \approx 0.0012$.

Three results follow, and they have the same shape.

\textbf{An unreported implementation choice moves the result by 45.9 points.}
Holding the backbone, the benchmark, the method, the number of layers removed
(four) and the compute saved ($-11\%$) all fixed inside one cell, and changing
only which layers were \emph{eligible} to be removed, success moves from
$+15.6$ to $-30.4$ against the same baseline. The ranking criterion is
identical in both; what differs is the candidate window and the spacing
constraint, values that the papers we surveyed generally do not report. Run in
all five cells, the same contrast spans 2.1 to 50.4 points --- so the
sensitivity to layer choice is itself a property of the configuration, not of
the method.

\textbf{A gain attributed to compression does not come from compression.}
Sweeping how much of the observation foveation keeps, the best setting is the
one that discards nothing: $+30.2$ points at \texttt{keep}$=100\%$ versus
$+18.8$ at $20\%$ and $+4.2$ at $10\%$. The intervention also changes measured
compute by between $-3.1\%$ and $+2.7\%$ across the five cells, because the
image size and the token count are unchanged. The effect is real and traceable
--- the log-polar round trip degrades the periphery irreversibly even when
nothing is subsampled --- but it is an input transformation, not an efficiency
technique.

\textbf{Changing the backbone moves results more than changing the
benchmark.} One benchmark-axis comparison passes correction; six
backbone-axis comparisons do, five of them at $p < 10^{-4}$. The cleanest case
is a sign reversal: on the same 135 Fractal episodes with the same
\texttt{depth prune 4}, OpenVLA gains $+15.6$ and SpatialVLA loses $-17.8$,
each significant on its own ($p = 4.3 \times 10^{-7}$ for the interaction).
The widest is \texttt{action repeat 2} on Bridge, where SpatialVLA gains
$+12.5$ and UniVLA loses $-69.8$ --- 82 points apart
($p = 1.7 \times 10^{-13}$) --- while the compute saved is identical by
construction.

This is not only visible in our own grid. Splitting the Google Robot results
that prior work already publishes, \texttt{pick coke can} rises or holds while
\texttt{move near} falls in all twelve configurations reported across two
author groups and three method families~\cite{efficientvla, vlacache, fastv}.
None of those papers discusses the split, because all of them report only the
four-task average.

Finally, we asked what breaks rather than how much. In the cell with the
largest drop we classified failures inside a single task and found that the
policy's ability to identify the commanded object was statistically unchanged
($1 \to 4$ episodes, $p = 0.375$), while episodes in which it never moved any
object at all went from $0$ to $12$ ($p = 0.0005$). Removing layers damages
execution before it damages grounding. We note that this contradicts the
hypothesis we designed the measurement to test.

\subsection*{Contributions}

\begin{itemize}
  \item A uniform re-measurement of three training-free VLA interventions over
        a $3 \times 2$ backbone-benchmark grid (five filled cells, eight
        conditions each, 7{,}198 episodes), with per-episode records released.
  \item Evidence that reported effects are frequently properties of the
        configuration rather than of the method: an unreported eligibility
        flag moves one result by 45.9 points, and the same intervention
        reverses sign across backbones with both directions individually
        significant.
  \item Evidence that compute saved does not predict success change. The
        intervention with the largest effect in our grid saves approximately
        no compute, and its benefit is maximised when it discards nothing.
  \item A measurement procedure for this class of claim --- episode-level
        pairing, an explicit determinism check, and a uniformity rule for the
        grid --- together with an account of what current reporting practice
        omits that makes such claims unreproducible.
\end{itemize}

We do not propose a new efficiency method, and we do not claim that any of the
three interventions fails in general. Our baselines reproduce published
numbers in four of five cells and fall 4.2 points short in the fifth, which is
inconsistent with a systematic setup error. What we claim is narrower and, we
think, more useful: the evidence currently offered for these methods does not
distinguish a property of the method from a property of the one configuration
it was measured in.
